%% kapa: exam=1 round=35 year=2024 period=2 form=A subject=appraisal_law subject_ko=감정평가관계법규 seq=34 number=34
%% source: Q-Net 감정평가사 기출문제 (cst00309) artlSeq=5213229 file=raw/1008/5213229/2교시 문제지 원본.pdf
%% conversion: kordoc 4.24.51 -> tools/gen1.py (auto); review=proofread-sonnet; flags=none
\begin{question}{34}
부동산등기법령상 등기사무에 관한 설명으로 옳지 않은 것은?
\begin{choices}
  \item 등기관은 접수번호의 순서에 따라 등기사무를 처리하여야 한다.
  \item 등기관은 등기사무를 처리한 때에는 등기사무를 처리한 등기관이 누구인지 알 수 없도록 조치하여야 한다.
  \item 토지등기부와 건물등기부는 영구히 보존하여야 한다.
  \item 등기부를 편성할 때 1동의 건물을 구분한 건물에 있어서는 1동의 건물에 속하는 전부에 대하여 1개의 등기기록을 사용한다.
  \item 폐쇄한 등기기록의 열람 청구는 관할 등기소가 아닌 등기소에 대하여도 할 수 있다.
\end{choices}
\end{question}
\answer{2}
